% Transcription — fonds Alexandre Grothendieck, shelfmark 72, batch 5 (pages 81-100).
% Established from the facsimile archives/batches/72/batch-05.pdf (from a copy
% of 72.pdf fetched through the project's relay,
% https://grothendieck-relay.onrender.com/source/72.pdf), per the skill
% transcribe-grothendieck. The raw scan's cover sheet is removed from the
% batch file: PDF page k is archivists' page 80+k, and the pencilled numbers
% bottom-left (« 81 » ... « 100 ») agree. The folder is 112 pages in six
% batches (1-20, 21-40, 41-60, 61-80, 81-100, 101-112), transcribed in parallel.
%
% Pass: Opus 5.5 (claude-opus-5-5), 2026-09-27 — first pass, unchecked against the pages by a human.
%
% Pages without a \page marker: 83 and 95. They are tinted
% wrapper sheets, otherwise blank, each inscribed in ink at the top right
% with a title for the leaves that follow: they take no \page marker and
% their words become \section headings, each with a \note. Page 83: « Groupes
% engendrés par des pseudo-réflexions projectives et relations avec les
% réalisations géom. des [struck word] hyperpolyèdres. » Page 95: « Groupes
% engendrés par des pseudo-réflexions vectorielles... » (first word doubtful
% there, read from the parallel page 83).
%
% Pages summarised: none. All eighteen other pages are his working notes,
% transcribed in full. Page 92 is cancelled by two long crossing diagonals
% and page 93's upper half by five oblique strokes; both are transcribed as
% far as they read, the cancellation said once. Page 86 is a sheet of
% formulas without prose, given in the order of the page.
%
% His pagination: none on these leaves. His numbered runs: formulas (1)-(4)
% on page 84; (5) (circled) and (F) on page 85; formulas (1)-(3) on page 91;
% the Théorème of page 96 numbers its formulas (1)-(16), continued through
% the proof on pages 97-100 to (20) and in Corollaire 2 (page 100) to (25').
% Page 92 cites (5), (8), (9) and F_n, F_{n-1}, which match no numbering on
% page 91: the order of the leaves around 91-93 is not certain.
%
% Arithmetic checked: page 93, tau_i^2 = 1 iff <a_i, a'_i> = lambda_i - 1 =
% -2 iff lambda_i = -1, correct. Slips of his flagged in \note, never
% corrected: page 81 writes h_{i+1} ∉ H_i where the inclusion Z_i ⊂ Z'_i
% seems to want ∈; page 85 writes s_{i+1} ∈ X_i and τ_i s_i ∈ X_i where the
% argument needs ∉; page 93 has two disagreeing values for <a_{i+1}, a'_i>
% and an undefined α_i; page 96, formula (14) reads E_i = k(s_0, ..., s_{i-1})
% where page 98 gives E_{n+1} = k(s_0, ..., s_{n+1}); page 100 drops a term
% in τ_i E_i (harmless for the conclusion).
%
% What the batch is about. Projective and linear realisations of regular
% « cartes » / quasi-polyhedra, through groups generated by pseudo-
% reflections. Pages 81-82 end a discussion (from batch 4) of subspaces
% invariant under G for a flag h_i, H_i: invariant Z lies between Z_i and
% Z'_i, conditions beta_i = 0; irreducibility iff all beta_i ≠ 0, i.e. the
% canonical bilinear form f_V non-degenerate. Pages 84-90 (after the cover of
% page 83): a regular combinatorial n-carte C with automorphism group G, its
% repères R a G-torsor, generators tau_i^(r); a geometric realisation in a
% projective bundle X is a hom G -> Aut X with a flag X_0 ⊂ ... ⊂ X_n,
% tau_i X_j = X_j (i ≠ j), tau_i reflections; a proposition of equivalent
% non-degeneracy conditions a) a') b) c) d) e) f) (s_{i+1} = tau_i s_i
% projectively independent; centres h_i, cocentres H_i with h_i ∉ H_{i+1});
% the lattice c, C, d, D, Y_i, Z_i (page 86); then the flags adapted to
% (tau_i) form a torsor under the automorphism group of (D_0, c, d), which is
% G_m when c ≠ d and G_a when c = d (pages 87-90). Pages 91-94: « coordonnées
% universelles », tau_i = id + a'_i ⊗ a_i, a_i = s_{i+1} - s_i, a'_i with
% parameters lambda_i, mu_i; the 2x2 matrices <a_i, a'_j>; a lemma lifting
% pseudo-reflections of P(E) to E; a list of things to « dégager ». Pages
% 96-100 (after the cover of page 95): Théorème, for pseudo-reflections tau_i
% of a k-module E and s_0 ∈ E, with s_i, E_i, lambda_i: conditions A (tau_i
% commute when non-consecutive, tau_i s_0 = s_0, (s_i) free) equivalent to
% B (<a_i, a'_j> = 0, <s_0, a'_i> = 0, (s_0, a_0, ..., a_n) free, lambda_n
% ∈ k*), with proof by induction on n, cqfd, and Corollaire 2 on
% h_i = k a_i, H_i = Ker a'_i.
%
% Notation, fixed once. His underlined Aut is \underline{\mathrm{Aut}} (the
% automorphism scheme); his script E is \mathcal{E}; « Lin » is
% \mathrm{Lin}; ∨ and ∧ are \vee, \wedge (join and meet); P with a check is
% \check{\mathbb{P}}. « Gram_i(X) » on pages 84 and 89 is kept as read
% (a Grassmannian is expected).
%
% The dating is the folder's own, copied verbatim from the catalogue; no
% group sentence.
\documentclass[11pt,a4paper]{article}
\input{../preamble/grothendieck}

\folder{72}
\batch{5}
\pages{81}{100}
\dating{[à partir de 1978]}
\watermark{Édition de démonstration}
\foldertitle{Polyèdres réguliers et hyperpolyèdres réguliers (Calculs en dimension quelconque) : notes manuscrites (s.d.).}

\begin{document}

\page{81}
\note{la page commence en cours de raisonnement.}
donc
\[
\underbrace{\mathrm{Lin}(h_{i+1}, \dots, h_n)}_{Z_i \ \mathrm{dim}\ n-i-1} \subset Z \subset \underbrace{H_0 \cap \dots \cap H_i}_{Z'_i \ \mathrm{dim}\ n-i} .
\]
\note{« $\mathrm{Lin}$ » est lu d'après la page 85, où le mot est écrit en clair ; ici il est abrégé. L'indice de dimension sous $Z_i$ se lit $n-i-1$, le dernier signe peu net.}

\struck{Cela implique donc $Z$}
et \uncertain{par} cette condition implique bien que $Z$ est stable par les $T_j$. Elle implique aussi \struck{$Z$}.
\[
Z_i \subset Z'_i ,
\]
ce qui équivaut aussi à
\[
h_{i+1} \notin H_i
\]
i.e.
\[
\beta_i = 0 .
\]
\note{le signe est bien $\notin$ sur la page ; l'inclusion $Z_i \subset Z'_i$ semblerait demander $h_{i+1} \in H_i$.}

\marginal{à droite, séparé du texte par un trait vertical : Réciproquement, pour $i$ fixé, conditions équiv.\ : a) $Z_i \subset Z'_i$ ; b) $Z_i$ inv.\ par $G$ ; c) \struck{$Z$} $Z'_i$ inv.\ par $G$ ; d) $\beta_i = 0$.}

Et pour une raison de dim., on doit avoir
\[
Z = Z_i \quad \text{ou} \quad Z = Z'_i ,
\]
l'inverse étant clair.

\emph{Remarque}. Conditions équivalentes, lorsque $\dim C \geq 2$ i.e.\ $n \geq 2$ : a) Il n'y a pas d'hyperplan de $C$ invariant par $G$ ; b) On a \struck{\ill{}} $\beta_0 \neq 0$, et \struck{\ill{}} \add{on n'a pas} $\bigl(\beta_1 = 0$ et $\underbrace{H_0 \cap H_1}_{\mathrm{dim}\ n-1} \subset C\bigr)$ \add{(i.e.\ $H_0 \cap H_1$ inv.\ \struck{\ill{}} par $G$)}.
\note{sur « $\beta_1 = 0$ » un $\neq$ a d'abord été écrit, puis sa barre biffée. Sous la parenthèse une accolade renvoie à :}
$\Updownarrow$ si $\beta_0$ ou $\mu_0 \neq 0$ :
\[
\beta_1 = 0 \quad \text{et} \quad \mu_1 \mu_0 - \beta_0 = 0
\]

\page{82}
\ill{} forment un drapeau.

c) Pour que $G$ opère sur $C$ \ill{} irréductiblement, il faut et suffit que tous les $\beta_i$ ($0 \leq i \leq n-1$) soient $\neq 0$, i.e.\ que la forme \struck{quadratique can.} \add{bil.\ can.} $f_V$ sur $V$ soit \struck{non dég.} non dégénérée (i.e.\ que $f_V$ soit lisse, et que \struck{p\ill{}} \add{si} $n$ impair, \struck{$2$ soit inv.\ sur $S$} \add{$k$ soit de car.\ $\neq 2$}).

\struck{c) Pour que}

\emph{Dém.} Dire que $Z$ est inv.\ par $T_i$ signifie que l'on a
\[
Z \ni h_i \quad \text{ou} \quad Z \subset H_i .
\]
\struck{Si $Z$ est} \add{\ill{}} invariant par tous les $T_i$, \struck{ou bien $Z \subset H_i$ pour tt $i$, donc $Z \subset \bigcap H_i = \{c\}$, donc $Z = \emptyset$ ou $Z = \{c\}$, ou bien $\exists\, i$} Notons que si $0 < i \leq n$\note{bornes peu nettes.}
\[
Z \subset H_i \Longrightarrow Z \subset H_{i-1} ,
\]
car sinon on aurait $h_{i-1} \in Z$, donc $h_{i-1} \notin H_i$, absurde. De même, si $h_j \in Z$, \struck{$h_j$} $0 \leq j < n$, donc $h_{j+1} \in Z$, car sinon on aurait $Z \subset H_{j+1}$, donc $h_j \notin H_{j+1}$, absurde.
\note{dans « $h_{i-1} \notin H_i$ » et « $h_j \notin H_{j+1}$ », la barre est tracée sur un $\in$ déjà écrit.}
Il en résulte que l'on a deux possibilités \struck{\ill{}} :

a) $Z \subset H_i$ $\forall\, 0 \leq i \leq n$, donc $Z \subset \bigcap_{0 \leq i \leq n} H_i = \{c\}$, donc $Z = \emptyset$ ou $Z = \{c\}$ --- qui est bien invariant ;

a') $Z \not\subset H_i$ $\forall\, 0 \leq i \leq n$, i.e.\ $Z \ni h_i$ $0 \leq i \leq n$, donc $Z \supset V(h_0 \dots h_n) = C$, donc $Z = C$ \struck{\ill{}} \add{$Z =$} $X$\note{la fin de la ligne est surchargée : après « $Z = C$ » un trait barré, puis « $X$ », avec « $Z =$ » au-dessus.} qui est bien invariant ;

b) On n'est pas dans a) ou a'), donc $\exists\, i$, $0 \leq i < n$, tel que $Z \subset H_j$ \add{si $0 \leq j \leq i$}, et $Z \not\subset H_j$ si $j \geq i+1$

\section{Groupes engendrés par des pseudo-réflexions projectives et relations avec les réalisations géom.\ des \struck{\ill{}} hyperpolyèdres.}
\note{ces mots, de sa main, sont écrits en haut à droite d'un feuillet teinté (page 83), par ailleurs blanc, qui précède les pages suivantes ; un mot biffé devant « hyperpolyèdres ».}

\page{84}
\struck{Soit $C$ $n$-maquette \add{régulière}, correspon}

Soit $C$ une $n$-carte régulière combinatoire, \struck{\ill{}} $R$ l'ens.\ de ses repères, qui est un torseur à gauche sous $G = \mathrm{Aut}\, C$, le commutant de $G$ étant formé des opérations des gr.\ cartographiques, engendré par $\sigma_0, \dots, \sigma_n$. Si $\underline{r} \in R$, il lui correspond des générateurs $\tau_i^{(r)}$ de $G$, définis par
\[
(1) \qquad \tau_i^{(r)} \underline{r} = \sigma_i \underline{r} .
\]
Une réalisation géométrique projective de $C$, dans \struck{\ill{}} $(n+1)$-fibré projectif $X$ sur $S$, revient à la donnée de

a) Opération $\varphi_X$ de $G$ sur $X$\note{l'indice de $\varphi$ est peu net ($X$ ou $r$).}
\[
(2) \qquad \varphi_X : G \to \mathrm{Aut}_S X = \Gamma\, \underline{\mathrm{Aut}}_S(X) ,
\]

b) drapeau \add{\ill{}} $\underline{r} = (X_0, X_1, \dots, X_n)$ \struck{$X$} dans $X$, satisfaisant
\[
(3) \qquad \tau_i X_j = X_j \quad \text{si } i \neq j
\]
\note{sous (3) une seconde ligne est biffée : « $\tau_i X_i \neq X_i$ ».}

(4) Les $\tau_i$ sont des réflexions de $P$\note{ou de $X$ ; la lettre est peu nette.}

[Ceci permet de \add{définir $R \to \mathrm{Drapeaux}(X)$ et des} applications $\mathrm{Dr}_H(C) \to \mathrm{Dr}_H(X)$.]

\emph{Prop.} Conditions équivalentes (si \struck{\ill{}} $S = \mathrm{Spec}(K)$ \add{le corps})\note{lecture incertaine de la parenthèse.} :

b) \struck{Aucune} Aucune des applications
\[
\Phi_i \to \mathrm{Gram}_i(X) \qquad 0 \leq i \leq n
\]
n'est constante, \struck{\ill{} fibre \ill{} générale}
\note{« $\mathrm{Gram}_i$ » est lu tel quel ; on attend une grassmannienne. Dans la marge gauche, une étiquette biffée et une flèche qui descend vers le groupe d'énoncés suivant.}

\struck{c) $\forall\, 0 \leq i \leq n$, on a $\tau_i X_i \neq X_i$.}

\note{les quatre énoncés suivants sont réunis par une accolade ; plusieurs lettres ont été récrites sur d'autres.}

d) Si $\underline{f_{i-1}} < \underline{f_i} \neq \underline{f'_i} < \underline{f_{i+1}}$ dans $\Phi_*(C)$, alors $f_i \neq f'_i$
\marginal{indépendant du repérage $r$ (une flèche vers d))}

a) L'image \add{$\Phi_0$} de $\underline{\Phi}_0$ engendre lin.\ $X$ (i.e.) \ill{}

\add{\ill{} de $X_0$ engendre lin.\ $X$}

a') L'intersection des hyperplans de $\Phi_n$ est vide

c) $\forall\, 0 \leq i \leq n$, $\tau_i X_i \neq X_i$

\page{85}
\marginal{en haut, encadré : \emph{NB} Dans f), les conditions $\alpha$) $\beta$) $\gamma$) ne font intervenir que $(\tau_i)$, non le choix de $s_0$, qui est précisé dans $\delta$) $\eta$) $\rho$)}

e) Posant $s_0 =$ unique pt de $X_0$,
\[
\begin{aligned}
s_1 &= \tau_0 s_0 \\
s_2 &= \tau_1 s_1 \\
&\cdots \\
s_{n+1} &= \tau_n s_n ,
\end{aligned}
\]
les $s_i$ sont \struck{\ill{}} projectivement indépendants, i.e.\ engendrent l'espace $X$.

\emph{NB} (5)\note{le chiffre 5 est cerclé.}
\[
s_i = \tau_{i-1} \cdots \tau_0 s_0 \qquad (1 \leq i \leq n+1)
\]
donc si $j > i$
\[
\tau_j s_i = (\tau_{i-1} \cdots \tau_0) \tau_j s_0 = s_i
\]
car $\tau_j s_0 = s_0$.

\marginal{dans la marge gauche, en oblique : Il y a une condition duale e'), avec \struck{$X_n$} $S_{n+1} = X_n$, $S_n = \tau_n X_n$, $S_{n-1} = \tau_{n-1} S_n$, $\dots$, puis une égalité aux indices empâtés « $S = \tau X_0$ », en exigeant $\bigcap_{0 \leq i \leq n+1} S_i = \emptyset$}

[f) \struck{Soit} Soit $h_i$ ($H_i$) le centre (\ill{}) des $\tau_i$. On a :
\note{la liste qui suit est écrite dans la colonne de droite, séparée par une accolade.}

$\alpha$) $h_i$ \uncertain{engendrent} $C$ ($(h_i)$ indép.)

$\beta$) $\bigcap H_i = c$ ($(H_i)$ indép.)\note{« $= c$ » est une lecture douteuse.}

$\gamma$) $h_i \notin H_{i+1}$ $(0 \leq i \leq n-1)$

$\delta$) $s_0 \in \underbrace{H_1 \cap \dots \cap H_n}_{\text{(droite } D)}$ (\emph{NB} c'est plus fort \ill{} que $\tau_i s_0 = s_0$ $1 \leq i \leq n$)

$\eta$) $s_0 \notin H_0$ ou encore $s_0 \neq c$ ($\{c\} = \bigcap H_i$) ou encore $\tau_0 s_0 \neq s_0$

$\rho$) $s_0 \notin C$ i.e.\ $s_0 \neq d$ où $\{d\} = D \cap C$.]
\note{les lettres grecques $\eta$ et $\rho$ sont lues d'après la forme ; la seconde est peut-être récrite.}

\emph{Schéma de dém.} On a les implications
\note{sur la page les flèches sont doubles ($\Rightarrow$, $\Leftrightarrow$) ; « tr » : trivial. Le schéma est reproduit avec ses sommets et ses flèches.}

\begin{tikzcd}
a) \arrow[r, Rightarrow, "\mathrm{tr}"] & b) \arrow[r, leftrightarrow, "\mathrm{tr}"] & c) \arrow[r, Rightarrow] \arrow[d, leftrightarrow, "\mathrm{tr}"] & e) \arrow[r, Rightarrow, "\mathrm{tr}"] \arrow[dr, Rightarrow, "\mathrm{tr}"'] & a) \\
a') \arrow[ur, Rightarrow, "\mathrm{tr}"'] & & d) & & a')
\end{tikzcd}

la seule implication non triviale étant
\note{la phrase s'arrête là ; au bout de la ligne, dans le cadre de la colonne de droite, est écrit « d) $\Rightarrow$ e) ».}

On prouve l'indép.\ des $(s_0, \dots, s_i)$ par \struck{\ill{}} réc.\ sur $i$ ; on \struck{Si elle est établie} \uncertain{prouvera en même temps}

\emph{Cor.} Sous les conditions précédentes, on a
\[
(F) \quad \Bigl\lbrace\ X_i = \mathrm{Lin}(s_0, \dots, s_i)
\]
($0 \leq i \leq n+1$, posant $X_{n+1} = X$).
\note{sous cette ligne, biffée : « $\tau_j s_i = s_i$ si $j > i$ ».}
Si \ill{} prouvé pour \struck{\ill{}} $i \leq i_0$, prouvons le pour $i = i_0 + 1$, i.e.
\[
\left\lbrace
\begin{aligned}
X_{i+1} &= \mathrm{Lin}(s_0, \dots, s_i, s_{i+1}) \quad (= \mathrm{Lin}(X_i, s_{i+1})) \\
\tau_j s_{i+1} &= s_{i+1} \quad (j > i+1)
\end{aligned}
\right.
\]
\note{sur la page : « si $j > i+1$ », puis une flèche vers « vérifié par (5) », et deux lignes biffées d'un gribouillis.}

Par déf.\ $s_{i+1} = \tau_i s_i$, or $\tau_i X_{i+1} = X_{i+1}$, car $s_i \in X_i \subset X_{i+1}$, donc $\underline{s_{i+1} \in X_{i+1}}$. \struck{Reste} \add{Il faut prouver que} Reste à prouver que $s_{i+1} \in X_i$, i.e.\ $\tau_i s_i \in X_i$\note{sic, deux fois $\in$ : le sens demande $\notin$.}, car \struck{$\tau_i$} $X_i = \mathrm{Lin}(X_{i-1}, s_i)$, donc
\[
\tau_i X_i = \mathrm{Lin}(\underbrace{\tau_i X_{i-1}}_{X_{i-1}}, \underline{\tau_i s_i})
\]
qui serait égal à $X_i$, contrairement à l'hyp.\ c).

\page{86}
\note{page de formules, sans phrases de liaison ; on les donne dans l'ordre de la page, de haut en bas et de gauche à droite.}

En haut à gauche, deux suites d'indices :
\[
a_0 \ \text{---}\ a_{i-1} \qquad \begin{array}{c} a_{\cdot}\ \overline{a_{n-1}\ a_n} \\ a'_{i+1} \ \text{---}\ a'_{n-1}\ a'_n \end{array}
\]
\note{l'indice du premier terme de la suite du haut, noté ici par un point, est illisible \ill{} ; une barre coiffe $a_{n-1}\,a_n$.}

\[
c = \bigwedge H_i \qquad C = \bigvee h_i \qquad d = \Bigl(\overbrace{\bigwedge_{1 \leq i \leq n} H_i}^{Z_1}\Bigr) \wedge C \qquad D = \underbrace{\bigvee_{0 \leq i \leq n-1} h_i}_{Y_{n-1}} \vee c
\]
\[
Y_{n-1} \cap H_n = Y_{n-2}
\]
\note{le signe $=$ est surchargé.}
\[
Y_{n-1} \cap (H_n \cap \dots \cap H_{i+1}) = Y_{i-1}
\]
Encadré :
\[
\left\lbrace
\begin{aligned}
&Z_1 \not\subset C \\
&c \notin Y_{n-1}
\end{aligned}
\right.
\]

Un schéma d'inclusions, les traits verticaux portant des dimensions :
\[
\begin{array}{ccc}
& \overbrace{\qquad}^{2} & \\
Y_{n-1} & \subset & X \\
{\scriptstyle n-i}\ \Big| & & \Big|\ {\scriptstyle n-i} \\
Y_{i-1} & \subset & Z_{i+1} = H_n \cap \dots \cap H_{i+1} \\
{\scriptstyle n}\ \Big| & {\scriptstyle 2} & \Big|\ {\scriptstyle i} \\
\emptyset & \subset & Z_1 \\
& {\scriptstyle 2} &
\end{array}
\]
\note{à droite du schéma : « \struck{\ill{}} $Y_0$ », un trait vers « $\emptyset$ — », puis « $h_0 \notin H_1$ », et plus bas, isolés, « $c$ $d$ ».}

Encadré :
\[
\begin{aligned}
X_n &\longmapsto X_i = X_n \wedge Z_{i+1} \\
X_0 &\longmapsto X_i = X_0 \vee Y_{i-1}
\end{aligned}
\]

\[
c = \bigwedge_0^n H_i \qquad d = \Bigl(\overbrace{\bigwedge_1^n H_i}^{Z_1}\Bigr) \wedge C = Z_1 \wedge C
\]
\[
C = \bigvee_0^n h_i \qquad D = \Bigl(\underbrace{\bigvee_0^{n-1} h_i}_{Y_{n-1}}\Bigr) \vee c = Y_{n-1} \vee c
\]
\[
\begin{aligned}
X'_i &= (Y_{i-1} \vee c) = Z_{i+1} \wedge C \\
X''_i &= Y_{i-1} \vee d = Z_{i+1} \wedge D
\end{aligned}
\]
\note{entre les deux membres de $X'_i$, une égalité biffée : \struck{$= C \wedge Z_{i+1}$}. Au-dessus du premier $=$ de $X''_i$, un petit « ? ».}
\[
Y_{n-1} \ \text{---}\ Z_1 \qquad\qquad Y_1 \quad Z_{n-1} = H_n \wedge H_{n-1}
\]

Correspondances, en colonne :
\[
\begin{array}{cc}
h_i & H_i \\
c & C \\
d & D
\end{array}
\]
\note{plus bas à gauche, biffé : « $\bigwedge h_i$ ».}

Encadré, sous une accolade :
\begin{itemize}
\item $h_i \subset H_j$ si $i, j$ non consécutifs
\item $h_i \not\subset H_{i+1}$
\item \struck{$(h_i)$ indép.}
\item $\bigvee h_i = C$ de codim $1$
\item $\bigwedge H_i = c$ de dim $0$ / $d$
\end{itemize}

Au pied de la page : $s_0$ longueur, $d$ infini\note{les deux lettres sont écrites au-dessus des deux mots.}.

\page{87}
\struck{\ill{}}
$h_0\ h_1\ \dots\ h_n$ sections du fibré projectif $X$

$H_0\ H_1\ H_n$ hyperplans ($X$ de rang $n+1$)
\note{le mot « hyperplans » est traversé d'un trait, peut-être une biffure.}
\marginal{encadré à droite : $(h_i, H_i)$ centre et cocentre de la pseudo-réflexion $\tau_i$ dans $X$ ; $\tau_i \tau_j = \tau_j \tau_i$ si $i, j$ non consécutifs, d'où \ill{}}

\[
\emptyset = X_{-1} \subset X_0 = \{s_0\} \subset X_1 \subset \dots \subset X_n \subset X = X_{n+1}
\]
\[
\tau_i X_j = X_j \quad (i \neq j) \qquad 0 \leq i, j \leq n
\]
\note{sur la page : « $\tau_i X_j = X_j$ si $i \neq j$ ».}
\[
\Longrightarrow \quad \underline{\Phi}_i \to \Phi_i
\]

A) Conditions équivalentes

\struck{a) $\forall\, 0 \leq i \leq n$, $\Phi_i$ non réduit à un pt}
\note{cette première rédaction de a) et b) est barrée en zigzag.}

a) $\Phi_0$ engendre lin.\ $X$ en chaque fibre

a') l'intersection des $Z \in \Phi_n$ est vide (\struck{\ill{}} en chaque fibre)

b) $\forall\, 0 \leq i \leq n$, $\Phi_i$ non réduit à un pt

c) $\forall\, 0 \leq i \leq n$ $\tau_i(X_i) \neq X_i$

d) Si $\underline{f_{i-1}} < \struck{\ill{}}\ \underline{f_i} < \underline{f_{i+1}}$ dans $\underline{\Phi}_*$, alors $\exists\, \underline{f'_i} \neq \underline{f_i}$ entre $\underline{f_{i-1}}, \underline{f_{i+1}}$, tels que $f_i \neq f'_i$
\note{un « $\exists$ » isolé est écrit à droite de c), et repris dans d).}

e) Posant $s_1 = \tau_0 s_0$, $s_2 = \tau_1 s_1 = \tau_1 \tau_0 s_0$, $\dots$, $s_{n+1} = \tau_n s_n = \tau_n \tau_{n-1} \cdots \tau_0 s_0$, les $(s_i)_{0 \leq i \leq n+1}$ engendrent \add{lin.} \struck{aff.\ $X$} \add{aff.\ $X_n$}
\note{la fin de la ligne est surchargée ; lecture incertaine des deux ajouts.}

e') Condition duale avec $S_n =$\ill{}, \struck{\ill{}} $S_{n-1} = \tau_n S_n$, $S_{n-2} = \tau_{n-1} S_{n-1} = \tau_{n-1} \tau_n S_n$, $\dots$, $S_{-1} = \tau_0 S_0 = \tau_0 \tau_1 \cdots \tau_n S_n$, on veut $\bigwedge S_i = \emptyset$.

\struck{f) \ill{} des conditions}

f) \ill{} l'une des \struck{\ill{}} conditions en ch.\ fibre
\note{les trois conditions qui suivent sont réunies par une accolade.}

$\alpha$) \struck{\ill{}} $(h_i)_{0 \leq i \leq n}$ libres $\Longleftrightarrow$ engendrent hyperplan $C$

$\alpha'$) $(H_i)_{0 \leq i \leq n}$ libres $\Longleftrightarrow$ $\bigcap_0 H_i = \{c\}$, $c$ section de $X$

$\beta$) $h_i \not\subset H_{i+1}$ en chaque fibre $0 \leq i \leq n-1$

\note{« en ch.\ fibres » est ajouté au-dessus de $\alpha$) ; la ligne $\alpha$) commence par un mot biffé.}

\struck{$\gamma$) posons $D_0 = H_n \cap \dots \cap H_1$ (droite), on a $D_0 \cap C = \{d\}$, $\dots$, $c \in D_0$, une section $D_0 \cap C = \{d\}$ $\dots$}
\note{bloc raturé en plusieurs couches ; on n'en donne que ce qui se lit.}

$\delta$) $s_0$ section de \struck{\ill{}}
\[
D_0^{**} = D_0 \smallsetminus \{c, d\} = D_0 - {}
\]
\ill{}
\note{après le dernier « $-$ », « \struck{sections} des 2 pts », et au-dessus « lieu des \uncertain{fixes} de $\tau_0 | D_0$ » ; lecture incertaine.}

\marginal{encadré à gauche, relié par une flèche à la liste : \emph{NB} $\alpha$) $\alpha'$) impliquent $\lbrace h_i \subset H_j$, $h_j \subset H_i$ si $0 \leq i, j \leq n$, $i, j$ non consécutifs (ce qui exprime \uncertain{encore} la commutation des $\tau_i$). f) implique $X_i = \bigvee(s_0, h_0, \dots, h_{i-1}) = \bigwedge(X_n, H_n, \dots, H_{i+1})$. $X_i$ et $X_j$ se déterminent mutuellement pour $i < j$ par $X_i = \bigwedge(X_j \cap H_j \cdots H_{i+1})$, $X_j = \bigvee(X_i, h_i, \dots, h_{j-1})$.}

B) Pour $(\tau_i)$ donné, pour qu'il existe loc.\ (\ill{}) un drapeau $(X_i)$ satisfaisant aux conditions équivalentes ci-dessus, il faut et suffit que $(\tau_i)$ satisfasse les conditions $(\alpha, \alpha', \beta)$ de (A f).

C) Sous les conditions de B (i.e.\ $\alpha$ $\alpha'$ $\beta$), l'application $(X_i) \mapsto s_0$ est une bijection de l'ens.\ \struck{des} \add{$\mathrm{Dr}(\tau_i)$} des drapeaux $(X_i)$ adaptés à $(\tau_i)$ et non dégénérés, sur l'ens.\ des sections de $D_0^{**}$. Plus gén., \struck{$\forall\, 0 \leq i_0 \leq n$,} \struck{bijection de $\mathrm{Dr}(\tau_*)$ sur l'ens.}
\[
(X_*) \longmapsto X_{i_0}
\]
est une bijection de $\mathrm{Dr}(\tau_*)$ sur l'ens.\ des sections de $D_{i_0}^{**}$, où $D_i^{**} = D_i \smallsetminus \{c_i, d_i\}$, $D_i =$ fibré en droites des $(i+1)$-\uncertain{espaces} compris entre $Y_{i-1} = \bigvee(h_0, \dots, h_{i-1})$ et $Z_{i+1} = \bigwedge(\struck{H_n}\ H_n, \dots, H_{i+1})$ (resp.\ de dim $i-1$ et $i+1$), $c_i$, $d_i$ étant les sections de $D_i$ définies resp.\ par \struck{$Y'_i = Y_{i-1} \vee Z_{i+1}$, $\dots$\ et $Z'_i = \bigvee(Z_i, Y_{i-1})$, $Y_{i-1} \subset Y'_i$, $Z'_i \subset Z_{i+1}$ [\emph{NB} $Z$ stable par $\tau_0 \dots \tau_i$, $\dots$\ $Z$ stable par $(\tau_i, \dots, \tau_n)$]}
\[
\begin{array}{ccc}
X'_i = (Y_{i-1} \vee c) & \text{et} & (Y_{i-1} \vee d) = X''_i \\
\Vert & & \Vert \\
(Z_{i+1} \wedge C) & & (Z_{i+1} \wedge D)
\end{array}
\]
où $D$ est défini \ill{} comme $D = \bigvee(h_0 \dots h_{n-1}, c)$.
\note{le passage biffé est raturé en zigzag ; on n'en donne que ce qui se lit.}

\page{88}
D) Sous les conditions de B), soit \struck{\ill{}} $\underline{g}$ le groupe des automorphismes de $(X, \struck{(\tau_*)})$, i.e.\ des $G_n$-automorphismes de $X$, où $G_n$ est le groupe engendré par les gén.\ $\tau_i$ ($0 \leq i \leq n$), liés par relations $\tau_i \tau_j = \tau_j \tau_i$ si $i, j$ non consécutifs. On a donc
\[
\underline{g} \longrightarrow \underline{\mathrm{Aut}}\bigl(D_0, (c) \cdot (d)\bigr) = \mathrm{Aut}(D_0, \Delta_0)
\]
où $\Delta_0$ est le diviseur sur $D_0$ défini par les sections $c$, $d$
\note{au-dessus de « $(c) \cdot (d)$ » : « id sur », et « produit de div.\ relatifs » ; au-dessus de $\Delta_0$ : « rel.\ sur » ; lectures incertaines.}
; ceci dit, cet homom.\ est un \emph{isom.} --- par suite
\[
D_0^{**} = D_0 \smallsetminus \Delta_0
\]
est un $\underline{g}$-torseur : $\exists$ loc.\ section (c'est la condition de B, C), et deux sections sont conjuguées par un unique \add{$G_n$-}automorphisme. Donc \struck{En} en vertu de C, item pour le schéma des \add{$(\tau_*)$-}drapeaux non dégénérés, qui est un torseur sous $\underline{g}$.

\marginal{à gauche : \emph{NB} $\underline{g}$ est lisse commutatif \add{\ill{}} de dim rel $1$, ses fibres sont $\mathbb{G}_m$ ($c \neq d$) (resp.\ \ill{} $\mathbb{G}_a$ ($c = d$)).}

\emph{Scholie}. Un système $(\tau_*, X_*)$ \add{non dégénéré} peut être considéré comme une \emph{rigidification} d'un système $(\tau_*)$ non dég., i.e.\ d'une représentation projective non dég.\ de $G_n$. La cat.\ (sur $S$ fixé) des $(\tau_*, \struck{X_*})$ non dég.\ est équivalente à celle des $(\tau_*, X_*)$ (i.e.\ des \emph{réal.\ géom.\ non dég.\ de $G_n$}) munies d'un torseur sous le groupe \add{$\underline{g}_0$} $\underline{\mathrm{Aut}}(D_0, \Delta_0)$ \add{rel.\ sur} correspondant à cette dernière$\dots$\ Utilisant la section $s_0$ distincte de $d$, \struck{comme} \add{comme pt 1} \struck{\ill{}}, $c$ comme $0$ et $d$ comme $\infty$ (si $c \neq d$ en tt fibre i.e.\ le centre $c$ n'est à l'$\infty$ sur aucune fibre) on trouve que ce groupe est $\mathbb{G}_m$, donc \struck{il s'agit de la} \add{dans ce cas} la cat.\ des $(\tau_*)$ non dég.\ avec centre partout à dist.\ finie équivaut à la cat.\ produit de la cat.\ des $(\tau_*, X_*)$ par celle des Mod.\ inv.\ sur $S$.

E) Soit $\Pi$ un quasi-polyèdre combinatoire \emph{régulier}, $G$ le groupe de ses automorphismes, $R$ l'ens.\ des repères, qui

\page{89}
est un torseur sous $G$. Le choix d'un $r \in R$ définit un \struck{épim.}
\[
G_n \xrightarrow{\ \varphi_r\ } G ,
\]
d'où un système $\tau^r_* = (\tau^r_i) = \varphi_r(\underline{\tau}_i)$, transformé en un syst.\ \emph{conjugué} (par $\mathrm{int}(g)$) si on change $r$ en $r' = gr$ ($g \in G$). Donnons-nous un hom.
\[
\varphi_G : G \longrightarrow \mathrm{Aut}\, X ,
\]
dire que $\tau^r_* = (\tau^r_i) = \varphi_G(\underline{\tau}^r_i) = (\varphi_G(\tau^r_i))$ est formé de \emph{quasi-réflexions}, et que la condition de non dég.\ de B est satisfaite, ne dépend pas du choix de $r$. Donnons-nous une repr.\ régulière
\[
\bigl(\varphi_G,\ \varphi_* = (\varphi_i)_{0 \leq i \leq n}\ (\varphi_i : \Phi_i(\Pi) \to \mathrm{Gram}_i(X))\bigr)
\]
de $\Pi$, d'où $r \mapsto \varphi_G(\underline{r}) = X_*$, dire que $(\tau^r_*, X_*)$ est non dég.\ ne dépend pas du choix de $\underline{r}$ --- \struck{intrinsèques} \add{d'ailleurs} les conditions \add{a) a') b)} (d) de A) sont intrinsèques. La première en vertu \ill{} équivaut à l'existence locale (Zar.) d'une « réal.\ géom.\ non dég.\ » de $\Pi$ associée à $\varphi_G$. \struck{Donc} \struck{Tel encore, on voit que le grp} Comme $G$ est simplement \emph{transitif} sur $R$, on voit d'ailleurs que tt objet géom.\ \ill{} associé intrinsèquement à une situation $\tau_*$ repr.\ $(\tau_*, X_*)$ non dég.\ [tel que $D_0$, et ses sections $c$, $d$] est défini intrinsèquement aussi en termes d'une $\varphi_G$ \add{repr.\ quasi-polyédrale n.d.} (repr.\ d'une réalisation géom.\ \add{$(\varphi_G, \varphi_*)$} de $\Pi$). Ceci dit, $\mathrm{Aut}\, \varphi_G$ opère sur $(D_0, c, d)$, et on trouve
\[
\underline{\mathrm{Aut}}(X, \varphi_G) \xrightarrow{\ \sim\ } \underline{\mathrm{Aut}}\bigl(D_0, (c) \cdot (d)\bigr)
\]
\note{au-dessus du second membre : « rel sur ».}
donc la cat.\ des $\varphi_G$ quasi-pol.\ non dég.\ équivaut à celle des couples \struck{\ill{}} $\bigl((\varphi_G, \varphi_*), T\bigr)$, \add{$(\varphi_G, \varphi_*)$ réal.\ géom.\ proj.\ n.\ dég.\ de $\Pi$}, $T$ un torseur sous $\underline{\mathrm{Aut}}(D_0, c, d)$ --- \ill{}
\note{la ligne continue au bas de la page, peu lisible ; on y devine « intéressant ».}

\page{90}
que ci-dessus pour $T$ comme l'une des structures quasi-pol.\ non dég.\ subordonnées à $\varphi_G$. Lorsque $c \neq d$ en chaque pt, $T$ s'interprète comme un torseur sous $\mathbb{G}_m$ (opérant par homothéties, cf.\ plus bas). Si $c = d$, alors $D_0^*$ \add{$= D_0 \smallsetminus c$} est un \struck{torseur fibré affine}, \struck{trivialisé par $s_0$} \struck{un fibré vectoriel} de rang $1$, \add{trivialisé par $s_0$} et \ill{} $\simeq D_0^*$ le fibré \uncertain{vectoriel} ass.\ à ses translations$\dots$
\note{la fin de la page est surchargée ; la suite des corrections est incertaine. Le reste de la page est blanc.}

\page{91}
\emph{Coordonnées universelles} \add{($n \geq 1$)}. On prend $C = V(h_0, \dots, h_n)$ comme hyperplan à l'$\infty$, \struck{alors les $s_i$ ($i \leq$) dans} $E = X \smallsetminus C$ comme fibré affine de rg $n+1$, les $s_i$ \uncertain{étant} donc des sections de $E$ formant une base affine \struck{et donc les $s_1 - s_0, \dots$ forment une base de l'espace vectoriel des translations} \struck{\ill{}. On choisit comme base $s_0, s_1, \dots, s_n$}.

$\mathcal{E}$ \emph{vectoriel} \add{$\mathcal{E}$} ext.\ de $\underline{\mathcal{O}}_S$ par \struck{\ill{}} l'espace des tr.\ $V$ de $E$ \add{enveloppe} vect.\ de $E$ s'identifie donc à $\underline{\mathcal{O}}_S^{n+2}$, en prenant \struck{\ill{}} $(\lambda_0 = 1, \lambda_1, \dots, \lambda_{n+1})$ correspondants
\note{la fin de la phrase est raturée en plusieurs couches.}
\[
\struck{s_0 + \lambda_1 (s_1 - s_0) + \lambda_2 (s_2 - s_0) + \dots + \lambda_n (s_n - s_0)}
\]
\note{ligne biffée, avec sous les différences les noms $e_0$, $e_1$, $\dots$, $e_n$.}
\[
\mathcal{E} = k s_0 \oplus k s_1 \oplus \dots \oplus k s_{n+1} \xrightarrow{\ \pi\ } \underline{\mathcal{O}}_S , \qquad \sum \lambda_i s_i \longmapsto \sum \lambda_i
\]

$\tau_i$ \add{ps.-}réflexions de $\mathcal{E}$ satisfaisant \add{$0 \leq i \leq n$}
\note{à gauche, biffés : « $\tau_i \tau_j = \tau_j \tau_i$ si $i \neq j$ » et « $\tau_i s_i =$ ». Les numéros (1), (2), (3) sont cerclés ; une flèche marquée « conditions imposées » pointe vers le groupe qu'ils forment.}
\[
\tau_i = \mathrm{id} + a'_i \otimes a_i
\]
\begin{itemize}
\item[(1)] $\langle a_i, a'_j \rangle = 0$ si $0 \leq i < j \leq n$, \struck{$i, j$ non cons.}\ \struck{$j \neq i+1$} \add{[$\langle a_j, a'_i \rangle = 0$ si $0 \leq i \leq j \leq n$, $j \neq i+1$]}
\item[(3)] \struck{$\langle s_0, a'_j \rangle = 0$ si $j \geq 1$}
\item[] \struck{(3) $\langle s_0, a'_0 \rangle$, $\langle s_0, a'_1 \rangle$, $\dots$, $\langle a_{n-1}, a'_n \rangle$ inversibles}
\item[(2)] $s_{i+1} = \tau_i s_i$ \quad $0 \leq i \leq n$
\end{itemize}
\note{les indices de la condition entre crochets sont surchargés ; lecture incertaine.}

\[
a'_j = \lambda_{jj} s_j + \lambda_{j,j+1} s_{j+1} + \dots + \lambda_{j,n+1} s_{n+1}
\]
exprime (1) :
\[
(\tau_i s_i =)\ s_i + \lambda_{ii} a_i = s_{i+1} \qquad \struck{a_i =} \ \lambda_{ii} a_i = s_{i+1} - s_i
\]
ops $\lambda_{ii} = 1$, donc
\[
\boxed{a_i = s_{i+1} - s_i} \qquad \lambda_{ii} = 1 \quad \text{i.e.} \quad a'_i = s_i + \lambda_{i,i+1} s_{i+1} + \dots + \lambda_{i,n+1} s_{n+1}
\]
exprime (1) et (2)
\[
\left.
\begin{aligned}
&\lambda_{i,i+2} = \lambda_{i,i+3} = \dots = \lambda_{i,n+1} \quad (\ldots\ \mu_i) \\
&\lambda_{i,i+1} = \lambda_i
\end{aligned}
\right|
\]
exprime (3).
\note{« ops » : on peut supposer. Dans la parenthèse, un mot illisible \ill{} puis une lettre surchargée, lue $\mu_i$.}

\page{92}
\note{page barrée de deux longs traits obliques qui se croisent ; on la transcrit dans la mesure où elle se lit. Elle renvoie à des conditions (5), (8), (9) et à des sous-espaces $F_n$, $F_{n-1}$ qui ne figurent pas à la page 91 : la suite des feuillets n'est pas assurée ici.}

on peut supposer ces conditions satisfaites, \add{donc aussi (9) vrai.} \struck{\ill{}} \add{donc aussi (9)}. \struck{Prouvons} \add{Notons} que alors \struck{l'indépendance} des \struck{\ill{}} conditions $s_{n+1} \notin$ \uncertain{libre} mod $F_n$ équivaut à la condition \struck{libre mod $F_n$}\ (5), \struck{Maintenant on sait déjà que} \add{Mais on sait que} \struck{$a'_n | F_{n-1} = 0$, donc} équivaut à $\langle s_n, a'_n \rangle$ \struck{\ill{}} \add{inv.}\ ?, et $\langle s_n, a'_n \rangle = \lambda_n$, OK. Il reste, pour compléter la récurrence, à établir (8) pour $i = n+1$ \add{i.e.} on a
\[
s_{n+1} \overset{\mathrm{déf}}{=} \tau_n s_n = s_n + \langle s_n, a'_n \rangle a_n \equiv \langle s_n, a'_n \rangle a_n \quad (F_n)
\]
\note{sous cette ligne, un calcul entièrement biffé et raturé : « $F_n$ $\dots$\ $\equiv \lambda_{n-1} a_{n-1} \langle s_{n-1}, a'_{n-1} \rangle a_{n-1} \equiv \langle s_{n-1}, a'_n \rangle a_{n-1}$ mod $F_{n-1}$ », « mod $F_n = \mathcal{O}_S(s_0, a_0, \dots, a_{n-1})$ ($F_{n-1}$) », puis, dans un cadre, « $s_n \equiv \lambda_{n-1} a_{n-1}$ », « $a'_n | F_{n-1} = 0$ », « $s_{n-1} \equiv \lambda_{n-1} a_{n-1}$ », « donc $\langle s_n, a'_n \rangle =$ \ill{} ». Seul reste non biffé, à droite :}
\[
\Vert \qquad \lambda_n a_n \ \text{par (9)} \qquad \text{OK.}
\]
\marginal{dans la marge gauche, en oblique : Les conditions de stabilité, évidentes$\dots$}

\emph{Corollaire}. Avec \struck{\ill{}} fibré vectoriel ($\mathcal{O}$-Mod.\ l.l.\ de t.f.) et conditions fibre par fibre.

\emph{Re Corollaire}. Conditions sur $E$ loc.\ libre de rang $n+2$. Conditions \struck{\ill{}} équivalentes sur le système $(s_0, \tau_0, \dots, \tau_n)$ \add{section, ps.-réflexions}, définissant les $s_i$.

\note{au pied de la page, écrit tête-bêche :}
\[
a'_i \otimes \underline{a_i} \qquad \underline{a'_{i+1}} \otimes a_{i+1} \qquad a_0 \ \text{---}\ a_n \quad a'_1
\]

\page{93}
\note{le haut de la page (jusqu'à « $i + 2 \leq j$ ») est barré de cinq longs traits obliques.}
\[
\underline{s_0\ a_0 \ \text{---}\ a_{i-1}}
\]
\[
\langle s_0, a'_0 \rangle = 1 \qquad \langle a_i, a'_{i+1} \rangle = \langle s_{i+1} - s_i, s'_{i+1} + \cdots \rangle = 1
\]
\note{le premier terme de la seconde égalité est récrit (d'abord $s_i$ ?).}
\[
\struck{\langle a'_i, a_j \rangle =} \quad \langle a_j, a'_i \rangle = \langle s_{j+1} - s_j, s'_i + \lambda_{i,i+1} s'_{i+1} + \dots + \lambda_{i,n+1} s'_{n+1} \rangle \qquad j \geq i+2
\]
\[
i \leq j - 2 \qquad i + 2 \leq j \qquad \lambda_{ij}
\]

Encadré :
\[
\begin{aligned}
\tau_i &= \mathrm{id} + a'_i \otimes a_i \\
a_i &= s_{i+1} - s_i \\
a'_i &= s'_i + \lambda_i s'_{i+1} + \mu_i (s'_{i+2} + \dots + s'_{n+1}) \qquad 0 \leq i \leq n-1 \\
a'_n &= s'_n + \lambda_n s'_{n+1}
\end{aligned}
\]
\note{la lettre devant la parenthèse est surchargée (d'abord « $+$ », puis une lettre lue $\mu_i$) ; à droite de l'encadré, « pour \ill{} » et « $\lambda_0 \dots \lambda_n$ », « $\mu_0 \dots \mu_{n-1}$ », en partie coupés.}

\struck{Quand a-t-on}
\[
\overbrace{Y_i}^{} \subset Z_{i+1} \qquad \bigvee(h_0 \dots h_{i-1}) \quad \bigwedge(H_n, \dots, H_{i+1})
\]
\note{ces deux lignes sont entourées d'un trait qui les relie à l'encadré.}
\[
\tau_i^2 = 1 \quad \text{i.e.} \quad \underbrace{\langle a_i, a'_i \rangle}_{\lambda_i - 1} = -2 \ ? \qquad \Longleftrightarrow \qquad \boxed{\lambda_i = -1}
\]
\note{le calcul est juste : $\langle s_{i+1} - s_i, a'_i \rangle = \lambda_i - 1$, qui vaut $-2$ pour $\lambda_i = -1$.}
\[
\tau_i = \mathrm{id} + a_i \otimes a'_i \qquad \struck{\ill{}} \qquad \tau_i^n = \mathrm{id} + P_n(\langle a_i, a'_i \rangle)
\]
\[
\tau_i \tau_{i+1} = \mathrm{id} + a'_i \otimes a_i + a'_{i+1} \otimes a_{i+1} + \underbrace{\langle a_{i+1}, a'_i \rangle}_{\alpha_i - \lambda_i} a'_{i+1} \otimes a_i
\]
\note{sous l'accolade, « $\alpha_i - \lambda_i$ » suivi d'un terme biffé ; la lettre $\alpha_i$ n'est pas définie sur la page (on attend $\mu_i$). Le terme croisé est bien $\langle a_{i+1}, a'_i \rangle\, a'_{i+1} \otimes a_i$ ; avec l'encadré, $\langle a_{i+1}, a'_i \rangle = \mu_i - \lambda_i$.}
\[
\left(
\begin{array}{ll}
\struck{\ill{}} & V_i \quad a_i, a_{i+1} \\
V_i^{\vee}\ \struck{\ill{}} & V'_i \quad a'_i, a'_{i+1}
\end{array}
\right)
\]
\[
\begin{pmatrix}
\lambda_i - 1 & 1 + \alpha_i - \lambda_i \\
1 & \lambda_{i+1} - 1
\end{pmatrix}
\qquad
\begin{pmatrix}
\langle a_i, a'_i \rangle & \langle a_{i+1}, a'_i \rangle \\
\langle a_i, a'_{i+1} \rangle & \langle a_{i+1}, a'_{i+1} \rangle
\end{pmatrix}
\]
\note{dans la première matrice les « $-1$ » sont écrits sur des termes raturés. Le coefficient $\langle a_{i+1}, a'_i \rangle$ y est écrit $1 + \alpha_i - \lambda_i$, et sous l'accolade du produit $\alpha_i - \lambda_i$ : les deux ne s'accordent pas, et la page ne tranche pas.}

En marge à gauche :
\[
\struck{\lambda_i \alpha_{i+1} + \lambda_i - \alpha_i} \qquad \struck{\ill{} \alpha_i + \lambda_i} \qquad t\ \begin{pmatrix} -2 & 2 + \alpha_i \\ 1 & -2 \end{pmatrix} \qquad 2 - \alpha_i
\]
\note{ce dernier bloc est barré d'un trait ondulé.}

\drawing{en haut à droite, au crayon puis à l'encre, un réseau de droites obliques se croisant en losanges, un sommet entouré d'un cadre ; aux nœuds, des couples d'indices peu lisibles, lus sous réserve « \uncertain{$0,0$} », « \uncertain{$1,0$} », « \uncertain{$0,1$} », certains raturés, et au milieu une inégalité entre deux d'entre eux.}

\page{94}
\emph{Lemme}. Soit $\mathcal{E}$ fibré vectoriel de rg $N \geq 3$, $X = \check{\mathbb{P}}(\mathcal{E})$.

a) Soit $\tau_X$ une \add{ps.-}réflexion de $X$. Alors $\exists !$ \add{ps.-}réflexion $\tau_{\mathcal{E}}$ de $\mathcal{E}$ qui lui donne naissance, propre \struck{\ill{}} si $\tau_X$ l'est.
\note{l'indice de $\tau$ dans « si $\tau_X$ l'est » est récrit.}

b) Soit $C$ un \struck{hyperplan} \add{sous-fibré proj.} de $X$, \struck{stable par $\tau_X$, de} $C = \check{\mathbb{P}}(V)$, \struck{supposons $\tau_X C = C$ i.e.\ $\tau_{\mathcal{E}} V = V$}. $\tau_X$ propre, centre $h \in C$, \struck{alors} alors \struck{considérons l'extension $\mathcal{E}$ de $\underline{\mathcal{O}}_S$ par $V$} $\tau_{\mathcal{E}}$ induit l'identité sur $\mathcal{E}/V$. Si $C$ est hyperplan et $\mathcal{E}$ l'enveloppe vectorielle can.\ de $E = X \smallsetminus C$, donc $\mathcal{E}/V \simeq \underline{\mathcal{O}}_S$ can., alors l'action de $\tau_X | E$ \struck{sur \ill{}} s'identifie à $\tau_{\mathcal{E}} | E$, où on identifie $E$ à l'image inverse de $1$ par $\mathcal{E} \to \underline{\mathcal{O}}_S$.

\emph{Dégager} : conditions de stabilité d'un s-fibré projectif par pseudo-réflexions

conditions de commutation de deux pseudo-réflexions.

\begin{itemize}
\item Ordre d'une pseudo-réflexion
\item Ordre du produit de deux pseudo-réflexions
\end{itemize}

\section{\uncertain{Groupes} engendrés par des pseudo-réflexions vectorielles$\dots$}
\note{ces mots, de sa main, sont écrits en haut à droite d'un feuillet teinté (page 95), par ailleurs blanc, qui précède les pages suivantes ; le premier mot est d'une lecture douteuse.}

\page{96}
\emph{Théorème}. Soient $k$ un anneau, $\mathcal{E}$ un $k$-module, $n \in \mathbb{N}$, \struck{\ill{}} $\tau_i$ ($0 \leq i \leq n$) des pseudo-réflexions dans $\mathcal{E}$
\[
(1) \qquad \tau_i = \mathrm{id} + a'_i \otimes a_i \qquad a'_i \in \check{\mathcal{E}}, \ a_i \in \mathcal{E} \quad (0 \leq i \leq n) .
\]
Soit $s_0 \in \mathcal{E}$, et posons \struck{Considérons}
\[
(2) \qquad s_1 = \tau_0 s_0, \ s_2 = \tau_1 s_1 \ (= \tau_1 \tau_0 s_0), \ \dots, \ s_{n+1} = \tau_n s_n \ (= \tau_n \tau_{n-1} \cdots \tau_0 s_0) ,
\]
\note{au-dessus : « $s_{i+1} = \tau_i s_i \dots$ ».}
\[
(3) \qquad \mathcal{E}_0 = k s_0, \ \mathcal{E}_1 = k(s_0, \struck{\ill{}}\, a_0), \ \struck{\ill{}} \dots, \ \mathcal{E}_i = k(s_0, \struck{\ill{}}\, a_0, \dots, \struck{\ill{}}\, a_{i-1}), \dots
\]
et enfin
\[
(4) \qquad \lambda_i = \langle s_0, a'_0 \rangle \langle a_0, a'_1 \rangle \cdots \langle a_{i-1}, a'_i \rangle \qquad (0 \leq i \leq n)
\]
\marginal{à droite, dans une accolade : $\mathcal{E}_{n+1} = k(s_0, a_0, \dots, a_n)$. \emph{NB} $\mathcal{E}_i$ engendré par $i+1$ \uncertain{éléments}.}
\struck{($\mathcal{E}_i$ engendré par $i+1$ \ill{}). On a alors}

1°) On a (5) $s_i \in \mathcal{E}_i$ pour \struck{\ill{}} \add{$0 \leq i \leq n+1$}, donc \struck{\ill{}} $k(s_0, \dots, s_i) \subset \mathcal{E}_i$ $(0 \leq i \leq n+1)$ \struck{\ill{}} (6) $s_{i+1} \equiv \lambda_i a_i$ mod $\mathcal{E}_{i}$
\note{dans (6), le membre de droite est récrit ; on lit $\lambda_i a_i$ modulo $\mathcal{E}_i$, sous réserve.}

2°) Conditions (A) et (B) \add{$(1 \leq i \leq n+1)$ équivalent à supposer $a'_i | \mathcal{E}_{i-1} = 0$ $(0 \leq i \leq n)$}
\note{l'ajout, relié par un trait à « (A) et (B) », est de lecture incertaine.}

A) \struck{Les sections $s_i$ ($0 \leq i \leq n+1$) sont lin.\ indépendantes} On a \struck{\ill{} $\mathcal{E}$ \ill{}}

(7) $\tau_i \tau_j = \tau_j \tau_i$ pour $0 \leq i, j \leq n$, $i, j$ non consécutifs

(8) $\tau_i s_0 = s_0$ pour $1 \leq i \leq n$

(9) $(s_i)_{0 \leq i \leq n+1}$ libre dans $\mathcal{E}$
\note{les numéros (7), (8), (9) sont récrits sur d'autres. À droite, une flèche vers (7)-(8) et, en oblique : « et il suffit que ce soit vrai modulo les \uncertain{scalaires} dans $\check{\mathbb{P}}(\mathcal{E})$ ».}

B) On a

(10) $\langle a_i, a'_j \rangle = 0$ \struck{\ill{}} pour $0 \leq i, j \leq n$, $i, j$ non consécutifs

(11) $\langle s_0, a'_i \rangle = 0$ pour $1 \leq i \leq n$

(12) $(s_0, \struck{\ill{}}\, a_0, \dots, a_n)$ libre dans $\mathcal{E}$.

(13) $\lambda_n \in k^*$

3°) De plus, ces conditions impliquent que
\[
(14) \qquad \mathcal{E}_i = k(s_0, \dots, s_{i-1}) \qquad 0 \leq i \leq n+1
\]
\note{d'après (3) et la \emph{NB}, $\mathcal{E}_i$ a $i+1$ générateurs $s_0, a_0, \dots, a_{i-1}$, et (5) donne $k(s_0, \dots, s_i) \subset \mathcal{E}_i$ ; on attendrait ici $k(s_0, \dots, s_i)$. L'indice $i-1$ est bien celui de la page.}
\struck{\ill{}}

(15) \struck{$\mathcal{E}_i$ stable par $\tau_0, \dots, \tau_{i-1}$ ($\dots i$)} $\tau_j(\mathcal{E}_i) = \mathcal{E}_i$ \quad $j \neq i$ \quad $0 \leq i, j \leq n$

(16) $\tau_i(\mathcal{E}_i) \neq \mathcal{E}_i$ \quad $0 \leq i \leq n$ \quad si $k \neq 0$

\page{97}
\emph{Dém.}

(1°) Prouvons la formule (5), qui se réduit à
\[
s_i \in \mathcal{E}_i \qquad 0 \leq i \leq n+1
\]
--- procédons par réc.\ sur $i$, la formule étant triviale pour $i = 0$. Et prouvons en même temps (6), (qui est vraie pour $i = 0$ si on pose $a_{-1} = s_0$, $\lambda_{-1} = 1$, $\mathcal{E}_{-1} = 0$). On a
\[
s_{i+1} = \tau_i s_i = \underset{\substack{\cap \\ \mathcal{E}_i}}{s_i} + \langle s_i, a'_i \rangle a_i \in \mathcal{E}_i + k a_i = \mathcal{E}_{i+1}
\]
\[
s_{i+1} \equiv \langle s_i, a'_i \rangle a_i \quad (\mathcal{E}_i)
\]
comme
\[
s_i \equiv \struck{\ill{}}\ \lambda_{i-1} a_{i-1} \quad (\mathcal{E}_{i-1})
\]
on aura
\[
\langle s_i, a'_i \rangle = \lambda_{i-1} \langle a_{i-1}, a'_i \rangle = \lambda_i ,
\]
d'où la formule (6) pour $i+1$ au lieu de $i$, si on suppose $a'_i | \mathcal{E}_{i-1} = 0$, OK.
\note{la dernière égalité utilise $a'_i | \mathcal{E}_{i-1} = 0$ et la définition (4) de $\lambda_i$ ; elle est juste sous cette hypothèse.}

\emph{Corollaire}. Soient
\[
\left\lbrace
\begin{aligned}
h_i &= k a_i \\
H_i &= \mathrm{Ker}\, a'_i
\end{aligned}
\right.
\qquad 0 \leq i \leq n \ \struck{\ill{}}
\]
les sous-modules associés à $\tau_i$, posons aussi $\lbrace a_{-1} = s_0$, $h_{-1} = k s_0$ \struck{\ill{}}.
Les conditions ci-dessous \struck{équivalentes} (A') et (B') sont équivalentes et impliquent la formule \uncertain{(14)} ($\mathcal{E}_i = k(s_0, \dots, s_{i-1})$)
\marginal{\emph{NB} $h_i$, $H_i$ \struck{sont} définis en termes des $\tau_i$, dès que \emph{$a_i$, $a'_i$} font libres.}

(A') On a

\struck{\ill{}} (15) $h_i \subset H_j$ si $-1 \leq i < j \leq n$, $j \neq i+1$

(15') \struck{\ill{}} $(s_i)_{0 \leq i \leq n+1}$ est libre (ou encore, \struck{\ill{}} [\struck{\ill{}}] les $h_i$ indépendants \struck{\ill{}} et $h_i \not\subset H_{i+1}$ en chaque fibre, pour $-1 \leq i \leq n-1$)
\note{l'ajout entre parenthèses, dans la marge de droite, est raturé en plusieurs couches ; lecture incertaine.}

(B') On a

(17) $\langle a_i, a'_j \rangle = 0$ pour $-1 \leq i < j \leq n$, $j \neq i+1$

(18) \struck{\ill{}} $(a_i)_{-1 \leq i \leq n}$ est libre

(19) $\lambda_n \in k^*$

Bien sûr, (17) n'est qu'une reformulation de (15), et il reste à prouver, en admettant (15) et (17) vrais, que (15') équivaut à (18) $+$ (19). \struck{Pour \ill{}} \struck{On \ill{}}. On peut supposer que \struck{\ill{}} ($S$ \add{i.e.\ $n \geq 0$}) \ill{}
\marginal{dans la marge gauche, en travers : $a_i \notin \mathcal{E}_{i-1}$ (lecture incertaine : l'appartenance est récrite).}

\page{98}
déjà prouvé pour $n-1$, et que $s_0 \dots s_n$ et $(a_{-1}, \dots, a_{n-1})$ sont lin.\ indép.\ et $(s_0, \dots, s_n)$ engendrent $\mathcal{E}_n$. \struck{On a donc \ill{}} \struck{Donc que $(s_0, \dots, s_n, s_{n+1})$ libre \ill{} $s_{n+1}$ libre mod $\mathcal{E}_n$ ; $s_{n+1} \equiv \lambda_n a_n$ mod $\mathcal{E}_n$, donc $s_{n+1}$ et $(s_0 \dots s_{n+1})$ libre \ill{}} \struck{$s_{n+1}$} or par (6)
\[
s_{n+1} \equiv \struck{\ill{}}\ \lambda_n a_n \quad (\mathcal{E}_n)
\]
donc la condition équivaut à $\alpha$) $a_n$ libre mod $\mathcal{E}_n$, i.e.\ (18), et $\beta$) \struck{\ill{}} la condition (19). \struck{\ill{}}
\note{les numéros (18), (19) sont récrits sur d'autres.}

Donc on a l'équivalence dite, et de plus ces conditions équiv.\ impliquent \struck{$\mathcal{E}_n = k(s_0, \dots, s_n)$} $\mathcal{E}_n$
\[
\mathcal{E}_{n+1} = \mathcal{E}_n + k s_{n+1} = k(s_0, \dots, s_n) + k s_{n+1} = k(s_0, \dots, s_{n+1}) ,
\]
ce qui achève la dém.\ du corollaire.
\note{ici $\mathcal{E}_{n+1} = k(s_0, \dots, s_{n+1})$, avec le même indice que $\mathcal{E}$ ; cela confirme que le $s_{i-1}$ de la formule (14), page 96, est un glissement d'indice.}

(2°)\note{le numéro est cerclé.} On procède par récurrence sur $n$, supposant ($n \geq 0$) que c'est prouvé déjà pour $n-1$. Donc on peut supposer que toutes les conditions, écrites dans A, B sont satisfaites pour le système $(s_0, \tau_0, \dots, \tau_{n-1})$, et qui a \struck{donc} \add{et de plus} \struck{$\mathcal{E}_n = k(s_0, \dots, s_n)$} (14) pour $0 \leq i \leq n$, et prouvons que alors les syst.\ de conditions
(A'') :
\begin{itemize}
\item[$\alpha$)] $\tau_n$ commute à $\tau_0 \dots \tau_{n-2}$
\item[$\beta$)] $\tau_n s_0 = s_0$
\item[$\gamma$)] $s_{n+1}$ libre mod $\mathcal{E}_n$
\end{itemize}
équivaut au syst.\ de conditions (B'') :
\begin{itemize}
\item[$\alpha$)] $\langle a_i, a'_n \rangle = 0$, $\langle a'_i, a_n \rangle = 0$ pour $0 \leq i \leq n-2$
\item[$\beta$)] $\langle s_0, a'_n \rangle = 0$
\item[$\gamma$)] $a_n$ libre mod $\mathcal{E}_n$
\item[$\delta$)] $\lambda_n \in k^*$
\end{itemize}
et que ces conditions impliquent
\[
\mathcal{E}_{n+1} = k(s_0, \dots, s_{n+1}) .
\]
(20)
\note{« (20) » est écrit seul au pied de la page. Dans $\alpha$) de (B''), le second crochet s'écrit $\langle a'_i, a_n \rangle$, les arguments dans l'ordre inverse de l'autre.}

\page{99}
Comme $s_{n+1} = \tau_n s_n = s_n + \langle s_n, a'_n \rangle a_n$, \add{la condition de ($\mathrm{A}''\gamma$)} $s_{n+1}$ libre mod $\mathcal{E}_n$, implique $a_n$ libre mod $\mathcal{E}_n$, donc $(s_0, a_0, \dots, a_n)$ libre.
\struck{En particulier, $(a_i, a_n)$ est libre si $0 \leq i \leq n-1$, donc $\tau_i$ commutant à $\tau_n$ si $\langle a'_i, a_n \rangle = \langle a_i, a'_n \rangle = 0$. De même $(s_0, a_n)$ étant libre, la condition $\tau_n s_0 = s_0$ \ill{} (\ill{}) \ill{} implique}
\note{passage encadré et barré de traits obliques.}
Donc on peut supposer, pour établir l'équivalence de (A'') et (B''), que $(s_0, a_0, \dots, a_n)$ est libre. Mais alors $(a_i, a_n)$ étant libre pour $0 \leq i \leq n-2$, la condition $\alpha$) de ($\mathrm{A}''\alpha$) s'interprète par ($\mathrm{B}''\alpha$), ($\mathrm{A}''\beta$) par ($\mathrm{B}''\beta$), il reste à montrer, \emph{moyennant les conditions précédentes supposées satisfaites}, que ($\mathrm{A}''\gamma$) équivaut à ($\mathrm{B}''\delta$), \struck{Or} \add{Or ($\mathrm{A}''\gamma$)}, par l'hypoth.\ faite de liberté de $a_n$ mod $\mathcal{E}_n$, équivaut à
\[
\langle s_n, a'_n \rangle \in k^* ,
\]
or par hyp.\ de récurrence et 1°) déjà prouvé, on a
\[
s_n \equiv \lambda_{n-1} a_{n-1} \quad \text{mod } \mathcal{E}_{n-1}
\]
et $a'_n | \mathcal{E}_{n-1} = 0$ par ($\mathrm{B}''\alpha, \beta$), donc
\[
\langle s_n, a'_n \rangle = \lambda_{n-1} \langle a_{n-1}, a'_n \rangle = \lambda_n
\]
donc on trouve bien la condition $\lambda_n \in k^*$ de ($\mathrm{B}''\delta$). Enfin il est évident que toutes ces conditions impliquent $\mathcal{E}_n + k s_{n+1} = \mathcal{E}_n + k a_n = \mathcal{E}_{n+1}$, et comme $\mathcal{E}_n = k(s_0, \dots, s_n)$ par hyp.\ de récurrence, on en déduit bien (20),

cqfd.

\page{100}
(3°)\note{le numéro est cerclé.} On a déjà prouvé (14), il faut prouver (15) (16).

\emph{Formule} (15) : Pour $j < i$ cela provient de $\tau_j = \mathrm{id} + a'_j \otimes a_j$, et $a_j \in \mathcal{E}_i$, pour $j > i$ cela provient de $a'_j | \mathcal{E}_i = 0$ \add{(moyennant les hyp.\ A) B))}, qui implique $\mathcal{E}_i \subset H_j$ $(\mathcal{E}^{\tau_j})$.
\note{sous $H_j$, un signe d'égalité vertical renvoie à « $\mathcal{E}^{\tau_j}$ ».}

\emph{Formule} 16. On a
\[
\mathcal{E}_i = \mathcal{E}_{i-1} + k a_{i-1} ,
\]
donc
\[
\begin{aligned}
\tau_i \mathcal{E}_i &= \underbrace{\tau_i \mathcal{E}_{i-1}}_{\mathcal{E}_{i-1}} + k \tau_i a_{i-1} \\
&= \mathcal{E}_{i-1} + k \underbrace{\langle a_{i-1}, a'_i \rangle}_{\in k^*} a_i = \mathcal{E}_{i-1} + k a_i
\end{aligned}
\]
\note{la première égalité omet le terme $k a_{i-1}$ de $\tau_i a_{i-1} = a_{i-1} + \langle a_{i-1}, a'_i \rangle a_i$ ; le résultat $\mathcal{E}_{i-1} + k a_i$ est celui de la page.}
\uncertain{libre} mod $\mathcal{E}_{i-1}$, cela et comme ($a_{i-1}$ et $a_i$) \struck{\ill{}} prouve notre assertion.
\note{l'ordre des deux dernières lignes sur la page est incertain.}

\emph{Corollaire} 2. Sous les conditions du \struck{\ill{}}) (\uncertain{équivalentes à} les conditions du Th.\ + 2° corollaire), \struck{\ill{}} on a, posant
\[
\left\lbrace
\begin{aligned}
\mathcal{E}^{\tau_i} &= \mathrm{Ker}\, a'_i \ (\overset{\mathrm{df}}{=} H_i) \\
k a_i &= \mathrm{Im}(\tau_i - \mathrm{id}) \ (\overset{\mathrm{df}}{=} h_i)
\end{aligned}
\right.
\qquad 0 \leq i \leq n
\]
\note{devant l'accolade, un premier essai biffé (« $H_i \overset{\mathrm{df}}{=}$ »).}
\begin{itemize}
\item[(21)] (l'accolade ci-dessus)
\item[(22)] $\mathcal{E}_i \subset H_j$ \struck{($\mathrm{Ker}\, a'_j$)} pour $0 \leq i < j \leq n$
\item[(23)] $\mathcal{E}_i \cap H_i = \mathcal{E}_{i-1}$,
\item[(24)] $\tau_i \tau_{i+1} \neq \tau_{i+1} \tau_i$, \emph{même modulo $k^*$} (si $k \neq 0$)
\item[(25)] \struck{$h_i \in H_i$} i.e.\ $\langle a_i, a'_{i+1} \rangle \in k^*$ (donc $\neq 0$ si $k \neq 0$) ($-1 \leq i \leq n$)
\item[(25')] $h_i \not\subset H_{i+1}$ si $k \neq 0$\note{sous $h_i$ et $H_{i+1}$, reliés par des signes $=$ verticaux : $k a_i$ et $\mathrm{Ker}\, a'_{i+1}$.}
\end{itemize}
\note{le numéro (21) est en face de l'accolade de définition ; (24), (25) sont récrits sur d'autres chiffres. Sous (25), biffé : « $k a_i$ — $\mathrm{Ker}\, a'_{i+1}$ ».}

La formule (22) est déjà prouvée, (21) 1° provient du fait que $a'_i$ \uncertain{est} libre, (21 2°) du fait que $a'_i : \mathcal{E} \to k$ est surjectif, car $\langle a_{i-1}, a'_i \rangle \in k^*$ (avec $a_{-1} = s_0$ !) grâce à (13) (4). \struck{Donc} La formule (23) provient de $\mathcal{E}_i = k(a_{-1}, \dots, a_{i-1})$, et que $H_i = \mathrm{Ker}\, a'_i$, $a'_i$ nul sur $a_j$ $(-1 \leq j \leq i-2)$, $\langle a_{i-1}, a'_i \rangle \in k^*$.

\end{document}
